% ═══════════════════════════════════════════════════════════════
% Section 8: Case Studies
% ═══════════════════════════════════════════════════════════════

\section{Case Studies}
\label{sec:case-studies}

To show the practical difference between CVSS-only and CVSS+NISS scoring, we score four
registry entries with both systems, one from each gap group plus a read-only case.

\textbf{Correction.} In v1.5, Cases 1--4 used identifiers whose registry entries describe
different techniques (QIF-T0015, T0034 and T0050 are directed-energy damage, working-memory
poisoning and hardware fault injection, not the P300, calibration and surveillance attacks
described), and their NISS vectors and scores did not match the registry. Some case scores
were also computed as an unweighted mean of six, which is not Equation~\ref{eq:niss}. This
revision replaces them with entries as stored in the live registry.

\textbf{Method and caveat.} NISS scores are recomputed from the stored vectors with
Equation~\ref{eq:niss} and match the stored values. CVSS v4.0 base scores were computed
from the stored base vectors with the open-source \texttt{cvss} Python library (v3.6) and
have not been checked against the FIRST.org calculator. The vectors were unchanged between
the snapshot and the live registry, but the snapshot file stored NISS scores that do not
reproduce from the equation (for example 3.4 for QIF-T0035 against 2.7). None of these
techniques has been executed against an implanted device; Case~5 is the only analysis of a
real vulnerability.

\begin{table}[H]
\centering
\caption{CVSS v4.0 vs.\ NISS for four registry entries (live registry).}
\label{tab:case-studies}
\small
\begin{tabularx}{\textwidth}{@{}l X l c c@{}}
\toprule
\textbf{Technique} & \textbf{Name} & \textbf{Status} & \textbf{CVSS} & \textbf{NISS} \\
\midrule
QIF-T0001 & Signal injection       & Theoretical  & 5.8 (Med)  & 6.1 (Med), PINS \\
QIF-T0035 & P300 interrogation     & Demonstrated & 6.9 (Med)  & 2.7 (Low) \\
QIF-T0058 & Calibration poisoning  & Demonstrated & 7.1 (High) & 6.0 (Med) \\
QIF-T0072 & Transducer inversion   & Demonstrated & 5.6 (Med)  & 2.0 (Low) \\
\bottomrule
\end{tabularx}
\end{table}

\subsection{Case 1: Signal Injection (QIF-T0001, Gap Group 3)}

CVSS rates this Medium (5.8): the vector requires physical access and attack preconditions,
with high integrity and availability impact on the system and high integrity impact on
subsequent systems. What it cannot express is the consequence for the person: current
delivered at the electrode-tissue boundary can cause tissue stress or involuntary activation.
\texttt{NISS:1.1/BI:H/CR:H/CD:H/CV:E/RV:P/NP:T} gives 6.1 and sets PINS because BI is High.
The technique is theoretical: its two listed sources are a position paper and a P300 probing
study, so the score describes a modeled impact, not an observed one.

\subsection{Case 2: P300 Interrogation (QIF-T0035, Read-Only)}

CVSS rates this Medium (6.9) for high confidentiality impact on the device and on subsequent
systems. The leaked data is the person's own evoked responses, which reveal recognition of
stimuli such as PINs, faces or places~\cite{martinovic2012feasibility}.
\texttt{NISS:1.1/BI:N/CR:H/CD:H/CV:E/RV:F/NP:N} gives
$(0 + 0.5 \cdot 6.7 + 0.5 \cdot 6.7 + 6.7 + 0 + 0)/5 = 2.68 \to 2.7$. The low score is
a direct consequence of the read-only ceiling (Section~\ref{sec:niss-limits}); the CR and CV
values, not the total, carry the information. The registry's CD:H for a read attack is a
scoring choice that warrants review.

\subsection{Case 3: Calibration Poisoning (QIF-T0058, Gap Group 2)}

CVSS rates this High (7.1) for integrity impact during calibration. Corrupted calibration
data makes the decoder learn wrong mappings between neural signals and intended actions,
and the user may adapt to the corrupted mapping. \texttt{NISS:1.1/BI:L/CR:H/CD:H/CV:I/RV:P/NP:T}
gives $(3.3 + 3.35 + 3.35 + 10 + 6.7 + 3.3)/5 = 6.0$; CV:I records that the user cannot
detect the corruption. The entry's source is Schroder et al.~\cite{schroder2025cyberrisks}; the
registry pairs it with signed calibration data and recalibration audits as controls.

\subsection{Case 4: Transducer Inversion (QIF-T0072, Gap Group 1)}

A speaker is re-tasked as a microphone. CVSS rates this Medium (5.6), and its
confidentiality metrics describe the impact well. NISS adds little:
\texttt{NISS:1.1/BI:N/CR:N/CD:N/CV:I/RV:F/NP:N} gives 2.0, and the dual-use class is
silicon-only. The case is included to show that NISS is not meant to raise every score.

\subsection{Case 5: Real-World Vulnerability Disclosure}

The \qif framework has been applied to real vulnerability research. During
systematic analysis of the BCI software ecosystem, we identified a multi-phase
exploit chain in an open-source library used in clinical and research BCI
pipelines.

The exploit chain demonstrates escalation from device-side vulnerabilities
(S2/S3) to potential impact on clinical decision-making through corrupted data.
CVSS scores the software vulnerabilities accurately; NISS records the downstream
risk to patients whose clinical care depends on the integrity of the data stream.

Responsible disclosure is in progress. Specific vulnerability details, including
affected software and CWE identifiers, will be published after coordinated
disclosure concludes.
